\documentclass[11pt]{article}
\usepackage[margin=1in]{geometry}
\usepackage{amsmath,amssymb}
\usepackage{graphicx}
\usepackage[hidelinks]{hyperref}

\newcommand{\B}{\mathcal B}
\newcommand{\K}{\mathcal K}
\newcommand{\wh}{\widehat}
\newcommand{\eps}{\varepsilon}

\title{A Literature-Implied Negative Answer to the BCAP-Free Calkin Embedding Question}
\author{FA Banach run \texttt{fa\_banach\_001}, agent \texttt{agent\_lane\_05}}
\date{2026-06-16}

\begin{document}
\maketitle

\section*{Status}

This is a lightweight literature-implied answer packet.  It records that
Remark 1 after Theorem 2.1 of Boedihardjo--Johnson is answered negatively once
the lane-5 norm-identification packet is combined with the Astala--Tylli
measure-of-noncompactness results.

The conclusion is not claimed as a new counterexample produced by the run.  The
run contribution is the identification of the Boedihardjo--Johnson
representation norm with the measure-of-noncompactness norm, which makes the
existing Astala--Tylli separation apply directly.

\section*{Source Question}

Boedihardjo and Johnson define, for an infinite-dimensional Banach space \(X\),
an ultrapower representation
\[
  f:\B(X)/\K(X)\longrightarrow \B(\wh X), \qquad \dot T\mapsto \wh T.
\]
Their Theorem 2.1 proves that if \(X\) has the \(\lambda\)-bounded compact
approximation property, then \(f\) is a norm-one \((\lambda+1)\)-embedding.
Immediately after the theorem they write:
\[
  \hbox{``We do not know whether Theorem 2.1 is true without the hypothesis
  that \(X\) has the BCAP.''}
\]

\begin{center}
\includegraphics[width=0.98\linewidth]{figures/open_problem_crop.png}
\end{center}

\section*{Identification}

For \(T\in\B(X)\), let
\[
  \gamma_X(T)=\inf\{\eps>0:T(B_X)\subset K+\eps B_X
  \hbox{ for some compact }K\subset X\}
\]
be the Hausdorff measure of noncompactness of \(T(B_X)\).  Equivalently,
\[
  \gamma_X(T)=\inf\{\|Q_E T\|: E\subset X,\ \dim E<\infty\},
\]
where \(Q_E:X\to X/E\) is the quotient map.  The equivalence follows because a
bounded subset of a finite-dimensional subspace is compact.

The previous lane-5 packet proved the exact norm formula
\[
  \|\wh T\|=\gamma_X(T) \qquad (T\in\B(X)).
\]
Thus the BCAP-free version of Theorem 2.1 is equivalent to asking whether, for
each Banach space \(X\), there is \(C_X<\infty\) such that
\[
  \|T\|_e\le C_X\,\gamma_X(T) \qquad (T\in\B(X)).
\]
The opposite inequality \(\gamma_X(T)\le \|T\|_e\) is automatic.

\section*{Literature Input}

Astala and Tylli studied the bounded compact approximation property through
measures of noncompactness \cite{AT}.  Later sources quote the relevant
consequences in exactly the form needed here:

\begin{itemize}
\item A later paper recalls that the Hausdorff measure of noncompactness
  induces a norm on \(\B(X,Y)/\K(X,Y)\), but that this norm need not be
  complete in general, citing Astala--Tylli, Theorem 2.5
  \cite[Subsection ``Measures of non-compactness'']{Kothe}.
\item The same source quotes Lebow--Schechter, supplemented by
  Astala--Tylli Theorems 2.2 and 2.3, for the estimate
  \[
    \|T\|_e\le \delta\,\gamma(T)
  \]
  whenever the range space has the \(\delta\)-compact approximation property
  \cite[Subsection ``Measures of non-compactness'']{Kothe}.
\item Skillicorn records that Zemanek noticed a result of Astala--Tylli
  providing the first example of a Calkin algebra with a non-unique norm
  \cite[Section 1]{Skillicorn}; in the surrounding discussion these
  non-equivalent norms are noted to fail completeness
  \cite[Section 1]{Skillicorn}.
\end{itemize}

Crossref metadata also records a later direct paper by D. Chen,
\emph{Characterizations of bounded compact approximation property by Calkin
representations}, Proc. Amer. Math. Soc. 150 (2022), 5397--5402,
DOI \texttt{10.1090/proc/16056}.  Its abstract says that it characterizes BCAP
via Calkin representations, and its reference list includes both
Boedihardjo--Johnson and Astala--Tylli.  The AMS PDF was blocked by Cloudflare
from this environment, so this packet does not quote Chen's theorem statement.

\section*{Consequence for the Boedihardjo--Johnson Question}

The Astala--Tylli/Skillicorn example gives a Banach space \(X\) for which the
measure-of-noncompactness norm \(\gamma_X\) on the Calkin algebra
\(\B(X)/\K(X)\) is not equivalent to the usual quotient norm
\(\|\cdot\|_e\).  Since always \(\gamma_X(T)\le\|T\|_e\), non-equivalence means
that no constant \(C_X\) can satisfy
\[
  \|T\|_e\le C_X\gamma_X(T)\qquad (T\in\B(X)).
\]
Equivalently, there is a sequence \(T_n\in\B(X)\) such that
\[
  \|T_n\|_e=1,\qquad \gamma_X(T_n)\longrightarrow 0.
\]
Using the lane-5 formula \(\|\wh T\|=\gamma_X(T)\), this becomes
\[
  \|T_n\|_e=1,\qquad \|\wh T_n\|\longrightarrow 0.
\]
Therefore the Boedihardjo--Johnson representation is not bounded below on the
Calkin algebra of this \(X\).  Hence Theorem 2.1 does not remain true for all
Banach spaces without the BCAP hypothesis.

\section*{Scope}

This packet answers the all-Banach-spaces form of Remark 1 negatively.  It does
not identify whether every non-BCAP space fails the representation embedding,
nor does it settle the isometry question in Remark 2 for \(L_p(0,1)\),
\(p\ne2\).

\section*{Verification Notes}

Human review should check one external identification: in the
Astala--Tylli/Skillicorn example, the non-equivalent Calkin algebra norm is
indeed the Hausdorff measure-of-noncompactness quotient norm \(\gamma\).  Once
that is confirmed, the argument above is immediate from the previous lane-5
norm-identification packet.

\begin{thebibliography}{9}

\bibitem{BJ}
M. T. Boedihardjo and W. B. Johnson,
\emph{On Mean Ergodic Convergence in the Calkin Algebras},
Proc. Amer. Math. Soc. 143 (2015), 2451--2457; arXiv:1301.0054.

\bibitem{AT}
K. Astala and H.-O. Tylli,
\emph{On the bounded compact approximation property and measures of
noncompactness},
J. Funct. Anal. 70 (1987), 388--401.

\bibitem{Skillicorn}
R. Skillicorn,
\emph{The uniqueness-of-norm problem for Calkin algebras},
arXiv:1507.08118.

\bibitem{Kothe}
P. Kolwicz, K. Leznicki, and L. Maligranda,
\emph{Essential norm of multiplication operators between Kothe spaces},
arXiv:2212.06723.

\end{thebibliography}

\end{document}
